
\section{Syntax and Semantics of \firstOrderLogic{}}

%Since \firstOrderLogic{} follows structured representations of a system, a \firstOrderLogic{} world consists in objects and relations between them.
The framework of \firstOrderLogic{} generalizes \propositionalLogic{} analogously to the generalization of factored system representations towards structured system representations \cite{russell_artificial_2021}.
This more expressive framework is designed to reason about relations and functions between objects.
To capture this framework by tensors we here first define the syntax and then investigate tensor representations of the semantics.

\subsection{Syntax}

A \firstOrderLogic{} theory consists in a finite set of constant symbols $\worlddomain$, a set of function symbols $\{\folfunctionof{\secatomenumerator}\wcols\secatomenumeratorin\}$, a set of predicate symbols $\{\folpredicateof{\atomenumerator}\wcols\atomenumeratorin\}$ and an arity map assigning the arity $\indorderof{\secatomenumerator}$ to the function $\folfunctionof{\secatomenumerator}$ and the arity $\indorderof{\atomenumerator}$ to the predicate $\folpredicateof{\atomenumerator}$.

\subsection{Tensor Representation of Semantics}

We here follow the model-theoretic semantics of \firstOrderLogic{}, where sets of possible worlds are considered.
% Database
We only treat in this work database semantics, where we assume that the domain $\worlddomain$ of each world has a one-to-one mapping to the constants of the theory and is therefore constant among the worlds.
Database semantics is central to combine the \semanticStructure{} with the \substitutionStructure{}.
Under this assumption, the dimension of the \substitutionStructure{} does not depend on the specific world, and we can combine both structures in a single tensor space to be defined in the next sections.

% Index interpretation of world domain
To each world there is a world domain $\worlddomain$ of objects, which we assume to be finite.
We exploit the set-encoding formalism discussed in more detail in \charef{cha:basisCalculus} and use bijective index interpretation maps
\begin{align*}
    \indexinterpretation \defcols [\inddim] \rightarrow \worlddomain \, .
\end{align*}
A so-called term variable $\indvariable$ takes states $\indindexin$, which represent objects
\begin{align*}
    \indexinterpretationat{\indindex} \in \worlddomain \, .
\end{align*}

% Relations
The relations between objects are described by $\indorder$-ary predicates $\folpredicate$.
Given a specific world $\worldindices$ the truth of relations is represented by boolean tensors
\begin{align*}
    \groundingof{\folpredicate} \defcols \symindstates\rightarrow\ozset \, .
\end{align*}
Given a tuple $\indindexlist\in\symindstates$ the boolean
\begin{align*}
    \groundingofat{\folpredicate}{\indexedindvariableof{0},\ldots,\indexedindvariableof{\indorder-1}} \in\ozset
\end{align*}
is called a grounding and encodes, whether the relation $\folpredicate$ is satisfied in the world $\worldindices$ for the objects $\indexinterpretationat{\indindexof{0}},\ldots,\indexinterpretationat{\indindexof{\indorder-1}}$.

% Functions
Functions in \firstOrderLogic{} are object-valued maps
\begin{align*}
    \groundingof{\folfunction}\defcols \worlddomain^{\inddim} \rightarrow \worlddomain \, .
\end{align*}
While relations are represented by their coordinate encodings, functions are represented by directed basis encodings
\begin{align*}
    \bencodingofat{\groundingof{\folfunction}}{\indvariableof{\folfunction},\shortindvariables}
\end{align*}
where to each object tuple $\indexinterpretationat{\indindexof{0}},\ldots,\indexinterpretationat{\indindexof{\indorder-1}}$ the vector $\bencodingofat{\groundingof{\folfunction}}{\indvariableof{\folfunction},\indexedshortindvariables}$ is the one-hot encoding of the object, that is
\begin{align*}
    \bencodingofat{\groundingof{\folfunction}}{\indvariableof{\folfunction},\indexedshortindvariables}
    = \onehotmapofat{\groundingofat{\folfunction}{\indexinterpretationat{\indindexof{0}},\ldots,\indexinterpretationat{\indindexof{\indorder-1}}}}{\indvariableof{\folfunction}} \, .
\end{align*}



%% NEW
\subsection{Worlds}

Worlds in First-Order Logics contain varying numbers of objects.
We use domain variables $\indvariable$ to enumerate the sets of objects.

Worlds are represented by three sets of tensors:
\begin{itemize}
    \item One-hot encodings of constants
    \item Coordinate encodings for each predicate
    \item Basis encoding of each function
\end{itemize}


%\subsection{Database semantics}

% Database Semantics
%\begin{align*}
%    \bigotimes_{\atomenumeratorin,\shortindindices\in[\inddim]^{\indorder}} \rr^2 \, .
%\end{align*}

\subsection{Two Tensor Structures}

In comparison with \propositionalLogic{}, \firstOrderLogic{} bears two natural tensor structures.
\begin{itemize}
    \item \textbf{\SemanticStructure{}:} As in \propositionalLogic{}, we enumerate possible worlds by a collection of variables $\catvariable$.
    This is a representation of the model-theoretic semantics, where subsets of worlds are represented by boolean tensors.
    \item \textbf{\SubstitutionStructure{}:}
    Different to \propositionalLogic{}, a formula can have variables and the evaluation of a formula on a world is represented by its grounding tensor.
    Given a world which contains objects in the domain $\worlddomain$, we can substitute variables by objects in the domain.
    In this way, each predicate with $\inddim$ variables is represented in that world as a boolean tensor of order $\inddim$.
\end{itemize}
% Russel reference
While the \semanticStructure{} appears already in factored representations of systems, the \substitutionStructure{} arises in more generality when treating structured representations of systems \cite{russell_artificial_2021}.





